% ============================================================
%  第八章
% ============================================================
\section{线算子：Wilson 与 \texorpdfstring{'t Hooft}{'t Hooft}}
\markboth{第八章\ 线算子：Wilson 与 't Hooft}{第八章\ 线算子：Wilson 与 't Hooft}

\noindent{\small\color{shuyiGray}对应底本：原 p.41--p.48。}\par\medskip

\subsection{什么是线算子}

\srcp{41}量子场论里的一个“线算子”，一般地说，是沿一条嵌入的一维流形 $\gamma$ 对理论所作的修改。$\gamma$ 可以是闭圈，也可以是端点在时空边界上的“开”曲线。

\begin{figure}[htbp]
\centering
\begin{tikzpicture}[scale=1.0]
  % 开曲线
  \fill[gray!12] (-4.6,-1.4) rectangle (-0.6,-1.05);
  \draw[gray] (-4.6,-1.05) -- (-0.6,-1.05);
  \node[below] at (-2.6,-1.42) {\scriptsize 时空边界};
  \draw[thick] (-3.9,-1.0) .. controls (-3.4,0.1) and (-3.0,0.5) .. (-2.6,0.0)
               .. controls (-2.2,-0.5) and (-1.8,0.3) .. (-1.3,-1.0);
  \node at (-3.75,0.55) {\small $\gamma$};
  \node[below] at (-3.9,-1.0) {\scriptsize 端点};
  \node[below] at (-1.3,-1.0) {\scriptsize 端点};
  % 闭圈
  \draw[thick] (1.9,0) circle (0.72);
  \draw[->,thick] (1.9+0.72*0.707,0.72*0.707) arc (45:80:0.72);
  \node at (1.9,1.0) {\small $\gamma'$};
\end{tikzpicture}
\caption{线算子：左边是端点在边界上的开曲线，右边是闭圈。沿 $\gamma$ 插入一个算子的意思，是要求理论在 $\gamma$ 的管状邻域里有指定的行为。}
\label{fig:line}
\end{figure}

\subsection{Wilson 算子}

\srcp{42}最简单的线算子是 Wilson 算子：它就是联络的全纯（holonomy），取规范群的某个表示 $R$。对闭圈，Wilson 算子定义为联络 $A$ 的全纯在表示 $R$ 下的迹：
\[
  W_R(\gamma)=\Tr_R\operatorname{Hol}_\gamma A
\]
（物理学家写作 $\Tr_R P\exp\oint_\gamma A$）。对开曲线，Wilson 算子需要稍长一点的讨论，我略去。注意：每个不可约表示 $R$ 都给出一个 Wilson 算子。

\subsection{\texorpdfstring{'t Hooft}{'t Hooft} 的原始定义}

\srcp{43}Wilson 算子在物理中被广泛用来探测夸克禁闭；大约在 1980 年，这促使 't Hooft 尝试定义 Wilson 算子的电—磁对偶。他给出的定义对他的应用来说是够用的，但事实证明，对几何 Langlands 而言，我们需要一个更精确的定义。原始定义实际上用的是我们在开头讨论过的朴素的磁荷定义。在四维流形 $M$ 上，一维流形 $\gamma$ 的余维是 $3$。't Hooft 设想用这样一个 $G$-丛 $E\to M\setminus\gamma$ 来做规范理论：它只在 $\gamma$ 的补集上有定义。由于四维中一维流形的“链环”（link）是二维球面 $S^2$，这种情况下就有一个局部不变量——它是 $S^2$ 上 $G$-丛的拓扑型，由 $\pi_1(G)$ 分类。't Hooft 原始定义的 't Hooft 算子，就由这个不变量分类。

\begin{yizhu}[disorder operator]
\textit{disorder operator}（无序算子）是相对 \textit{order operator}（有序算子）
而言的：后者由场的本地表达式定义（Wilson 圈即是），前者由“在奇点附近改变路径积分
的边界条件”定义，一般没有本地表达式。中文文献或译“无序算子”，本书随物理用法称
“线算子”，必要时括注英文。
\end{yizhu}


\begin{figure}[htbp]
\centering
\begin{tikzpicture}[scale=1.05]
  \shade[ball color=gray!6] (0,0) circle (1.0);
  \draw[thick] (0,0) circle (1.0);
  \draw[densely dashed] (-1.0,0) arc (180:360:1.0 and 0.34);
  \draw (-1.0,0) arc (180:0:1.0 and 0.34);
  \fill (0,0.26) circle (1.5pt);
  \node[right=3pt] at (0,0.26) {\scriptsize $p$};
  \node[above] at (0,1.1) {\small $S^2$};
  \draw[->,gray] (0.95,-0.7) -- (0.35,-0.25);
  \node[right,gray] at (0.95,-0.73) {\scriptsize $\gamma$ 的链环};
\end{tikzpicture}
\caption{四维中余维 $3$ 的线算子 $\gamma$：围绕它的链环是 $S^2$，其上的 $G$-丛拓扑型给出朴素磁荷。}
\label{fig:link}
\end{figure}

\subsection{Kapustin 的精确化}

\srcp{44}不过，Kapustin（2005）给出了一个更精确的定义。他的讨论与 't Hooft 所说不同之处在于：他对奇点附近发生的事情施加了更精确的条件，而不是仅仅指定情形的拓扑。

\srcp{45}想法是这样的：要求在场在 $\gamma$ 附近渐近于杨--米尔斯方程的一个指定的奇异解。所考虑的奇点如下。对余维 $3$ 奇点的局部模型，我们可以取 $\RR^3$ 原点处的行为。对 $G=U(1)$，荷 $1$ 的“狄拉克单极子”解是一个曲率为
\[
  F_0=\frac{1}{2}\star\dd\frac{1}{|\vec x|}
\]
的联络 $A_0$。对任何 $G$（Lie 代数为 $\mathfrak g$），取一个同态 $\rho:\mathfrak u(1)\to\mathfrak g$，并要求在 $\vec x=0$ 附近联络 $A$ 渐近于 $A_\rho=\rho(A_0)$。

\srcp{46}支在 $\gamma$ 上的、型为 $\rho$ 的 't Hooft 算子，由这样一条处方定义：“用沿 $\gamma$ 的这种奇点来做规范理论。”由于我们只关心 $A_\rho$ 在共轭意义下的类，就可以把 $\rho$ 看成到极大环面 $\mathfrak t$ 的同态：$\rho:\mathfrak u(1)\to\mathfrak t$。当然，我们也只关心 $\rho$ 在 Weyl 变换下的类。

\subsection{总结}

\srcp{47}总结一下：

（1）$G$ 规范理论中的 Wilson 算子由 $G$ 的表示分类，而表示又由最高权分类。一个最高权是一个同态 $\rho_{\mathrm{el}}:T\to U(1)$。

（2）$G$ 规范理论中的 't Hooft 算子由奇点分类，而奇点由同态 $\rho_{\mathrm{mag}}:U(1)\to T$ 标号。

\begin{table}[htbp]
\centering\small
\caption{电与磁在线算子上的对偶：表示与余权。}
\label{tab:line-ops}
\begin{tabular}{@{}lll@{}}
\toprule
 & Wilson 算子 & 't Hooft 算子 \\
\midrule
定义 & 沿 $\gamma$ 的全纯的迹 & 沿 $\gamma$ 的指定奇点 \\
标号 & $G$ 的表示 $R$ & 同态 $\rho_{\mathrm{mag}}:U(1)\to T$ \\
群论数据 & 最高权 $\rho_{\mathrm{el}}:T\to U(1)$ & 余权 $\rho_{\mathrm{mag}}:U(1)\to T$ \\
取值格 & $X^*(T)$ & $X_*(T)$ \\
对偶下 & $G$ 的 Wilson & $\gvee$ 的 't Hooft \\
\bottomrule
\end{tabular}
\end{table}

\subsection{对偶的陈述}

\srcp{48}现在我们可以陈述 Wilson 算子与 't Hooft 算子之间的对偶了：在规范群为 $G$ 或 $\gvee$ 的两个理论之间的 Montonen--Olive 对偶里，$G$ 的 Wilson 算子对应于 $\gvee$ 的 't Hooft 算子，反之亦然；所用的对应是 $\rho:T\to U(1)$ 与 $\rho^\vee:U(1)\to\Tvee$ 之间的自然对应，再加上 $G$ 与 $\gvee$ 有相同的 Weyl 群这一事实。这与 Atiyah 在 1977 年向我指出的那个原始的 GNO 对偶是同一回事；只不过，把它改述成线算子而不是粒子，就不必动用希格斯机制了。

\begin{pizhu}{余维 $3$：为什么链环是 $S^2$，为什么这是“局部”的}
在四维流形里，一维子流形 $\gamma$ 的管状邻域是 $\gamma\times D^3$，其边界是 $\gamma\times S^2$。要指定“沿 $\gamma$ 的奇点”，可以只指定一个横截球面 $S^2$ 上的行为；而 $S^2$ 上的 $G$-丛由 $\pi_1(G)$ 分类（把 $S^2$ 分成两个圆盘，转移函数在赤道上取值于 $G$）。

于是 't Hooft 算子天然是“局部的”：它只依赖 $\gamma$ 附近的数据，不依赖整体的真空。这与 Wilson 算子正好互补——后者是整体联络的全纯，依赖 $\gamma$ 所环绕的区域。两者放在一起，才构成一个完整的“电—磁”对。
\end{pizhu}

\begin{pizhu}{$\pi_1(G)$ 不够：Kapustin 到底改进了什么}
若 $G$ 单连通（如 $SU(N)$），$\pi_1(G)=0$，'t Hooft 的原始定义就什么也没给出：球面上的一切 $G$-丛都平凡。但物理上，$SU(N)$ 理论显然有丰富的 't Hooft 算子（例如把 $SU(2)$ 破缺到 $U(1)$ 后，磁荷可以取任意整数）。

Kapustin 的改进是：不只看拓扑，而是要求场在 $\gamma$ 附近\emph{渐近于一个指定的奇异解}。这个解由 $\rho:\mathfrak u(1)\to\mathfrak t$ 决定，而 $\rho$ 的共轭类就是余权。对 $SU(2)$，$\mathfrak t\cong\mathfrak u(1)$，$\rho$ 由整数 $n$ 标号；于是 't Hooft 算子有无穷多个，尽管 $\pi_1(SU(2))=0$。

这个区别正是几何 Langlands 所需要的：Hecke 修改的类型是余权，而不是 $\pi_1(G)$ 的元素。第 39 页的“类型由对偶群的表示标号”，在这一章变成了“\ 't Hooft 算子的型由 $\gvee$ 的表示标号”。
\end{pizhu}

\begin{pizhu}{为什么改述成线算子就不必希格斯化}
1977 年 Witten 的异议是：GNO 的匹配只在把 $G$ 破缺到 $T$ 之后才成立，而一个根本性的陈述不应依赖真空的选择。

线算子的表述绕开了这一点：Wilson 算子由表示 $R$ 定义，'t Hooft 算子由奇点型 $\rho$ 定义；两者都在\emph{任何}真空、\emph{任何}耦合下存在，不需要先选一个希格斯真空把 $G$ 破缺掉。对偶说的不是“两个破缺后的理论在低能下匹配”，而是“两个理论的算子代数在同构下匹配”。这就是 Witten 在 1977 年想要而当时没有的东西——它花了将近三十年。
\end{pizhu}

\begin{pizhu}{线算子对偶的群论内容}
把表 \ref{tab:line-ops} 的最后两行并排看：
\[
  X^*(T)\ \longleftrightarrow\ X_*(\Tvee),\qquad
  X_*(T)\ \longleftrightarrow\ X^*(\Tvee),
\]
这正是第 3 章的根数据对偶。于是“Wilson $\leftrightarrow$ 't Hooft”这件事，在群论层面就是“特征标格与余特征标格互换”。物理上它是一个关于耦合的断言（$e\leftrightarrow e^\vee$）；数学上它是一个关于格与表示的断言。两者能对上，是因为对偶群 $\gvee$ 的定义方式恰好让这两个格互换。

Witten 说“这与 1977 年的 GNO 对偶是同一回事”，指的就是这个：粒子层面的 GNO 对偶与算子层面的对偶，共享同一个根数据对偶，只是一个需要希格斯化、一个不需要。
\end{pizhu}
